\documentclass[conference]{IEEEtran}
\IEEEoverridecommandlockouts
% The preceding line is only needed to identify funding in the first footnote. If that is unneeded, please comment it out.
% Template version as of 6/27/2024

\usepackage{cite}
\usepackage{amsmath,amssymb,amsfonts}
\usepackage{algorithmic}
\usepackage{graphicx}
\usepackage{textcomp}
\usepackage{xcolor}
\usepackage{listings}
\usepackage[hidelinks]{hyperref}
\lstset{
    basicstyle=\ttfamily\footnotesize,
    breaklines=true,
    breakatwhitespace=true,
    columns=fullflexible,
    frame=single,
    keepspaces=true,
    showstringspaces=false
}
\def\BibTeX{{\rm B\kern-.05em{\sc i\kern-.025em b}\kern-.08em
    T\kern-.1667em\lower.7ex\hbox{E}\kern-.125emX}}

\begin{document}

\title{Reliable Data Transfer Protocol over UDP\\
{\footnotesize CS 5470 Computer Networks}
}

\author{
\IEEEauthorblockN{Joe Jimenez}
\IEEEauthorblockA{\textit{California State University, Los Angeles}\\
Los Angeles, California}
\and
\IEEEauthorblockN{Alan Lam}
\IEEEauthorblockA{\textit{California State University, Los Angeles}\\
Los Angeles, California}
}

\maketitle

\begin{abstract}
UDP is useful when an application needs simple communication with little
transport overhead, but it does not guarantee delivery, ordering, duplicate
protection, or protection from corrupted data. This project implements a
reliable data transfer protocol over UDP in Python. The protocol uses sequence
numbers, acknowledgments, CRC checksums, timeouts, retransmission, message
packetization, and a configurable sender window. The implementation was tested
with local UDP sockets, Linux \texttt{tc netem}, and project scripts to observe
the behavior under delay, packet loss, corruption, duplication, reordering, and
configured retry and timeout limits.
\end{abstract}

\begin{IEEEkeywords}
UDP, reliable data transfer, stop and wait, sliding window, checksum
\end{IEEEkeywords}

\section{Introduction}
User Datagram Protocol (UDP) gives applications a small and fast transport
interface, but it leaves reliability to the application. The sender does not
automatically know whether the receiver accepted the data, and many applications
that want UDP's low overhead still need ordered, uncorrupted delivery.

This project implements those reliability mechanisms on top of UDP. Each data
packet carries a sequence number, acknowledgment number, payload, and CRC
checksum. The sender transmits a packet and waits for an acknowledgment; if the
acknowledgment does not arrive before the timeout expires, the sender resends
the packet. The sender supports stop and wait as well as a configurable sending
window, so the same code can be tested with one or many outstanding packets.
The receiver validates each incoming checksum, acknowledges valid packets,
buffers in-window packets, and delivers data upward in sequence order.

The packet format is kept simple for testing:
\texttt{seq\_num\allowbreak|\allowbreak ack\_num\allowbreak|\allowbreak payload\allowbreak|\allowbreak checksum}.
This format is easy to print and inspect while verifying behavior.
Testing was done with local UDP sockets and with controlled network conditions
introduced by Linux \texttt{tc netem} and project scripts that inject delay,
packet loss, corruption, duplication, reordering, and configured retry and
timeout limits.
\section{Background and Motivation}
UDP does not provide reliable delivery. This is a real problem for any application that wants UDP's low overhead but still needs to know whether data arrived correctly. Currently an application cannot tell if a datagram was lost, corrupted, duplicated, or delivered out of order unless there is extra logic.

The textbook framing of this problem is reliable data transfer (RDT), developed
incrementally in Kurose and Ross Chapter 3.4 from a perfect channel up through
pipelined protocols such as Go-Back-N and Selective Repeat \cite{b3}. TCP is the
production example of a reliable transport protocol, and it uses sequence
numbers, acknowledgments, retransmission, and flow control to provide a
reliable byte stream \cite{b1}. The RDT progression studies these same
mechanisms in a smaller setting, so the behavior is easier to inspect.

The current state of practice is not limited to TCP. QUIC is a modern transport
protocol that runs over UDP and provides reliable streams, loss recovery, flow
control, and security as part of the transport design \cite{b2}. This project
does not try to reproduce TCP or QUIC. Instead, it implements the core reliability
ideas over UDP so the main protocol behavior can be tested directly.
\section{Protocol Design}

\subsection{Packet Format}
Each packet uses four fields: sequence number, acknowledgment number, payload,
and checksum. For debugging, the project uses a simple text format:
\texttt{seq\_num|ack\_num|payload|checksum}. This made it easier to print
packets and confirm that the sender and receiver parsed the same fields. The
tradeoff is that the format is not binary efficient and the payload cannot
contain the \texttt{|} character.

\subsection{Checksum}
For error detection the protocol uses a Cyclic Redundancy Check (CRC), as
described in Kurose and Ross \S6.2.3 \cite{b3}. CRC treats the packet header
and payload as a binary string $D$, appends $r$ zero bits, and divides the
result modulo 2 by a fixed generator $G$ of $r{+}1$ bits. The remainder $R$ is
the checksum that travels in the packet. The receiver runs the same division
on the received bits and compares the recomputed remainder against $R$; if
they differ, the packet is treated as corrupted and dropped.

This implementation uses the 9-bit generator \texttt{111010101}, so the
checksum is the 8-bit CRC remainder. An $r$-bit CRC catches every error burst
of length $r$ or fewer and detects longer bursts with high probability
\cite{b3}, which is why CRC is a stronger choice than the simple byte-sum
checksum the project used earlier (still available as a fallback in
\texttt{calculate\_checksum} when \texttt{use\_crc=False}). On the sender side
the checksum is computed over \texttt{seq\_num}, \texttt{ack\_num}, and
\texttt{payload} before the packet is sent; the receiver recomputes it the
same way on arrival.

\begin{lstlisting}[language=Python]
generator = '111010101'
data_str = f"{seq_num}|{ack_num}|{payload}"
data_as_binary_str = string_to_binary(data_str)
data_as_binary_str = data_as_binary_str + ('0' * (len(generator)-1))
remainder = divide_modulo_2(data_as_binary_str, generator)
return int(remainder, 2)
\end{lstlisting}

\subsection{Acknowledgments and Retransmission}
The protocol uses positive acknowledgments. After sending a packet, the sender
waits for an ACK that contains the acknowledged sequence number. If the ACK
does not arrive before the timeout expires, the sender retransmits the packet.
The sender also ignores malformed acknowledgments and acknowledgments with bad
checksums. This gives the sender a direct recovery mechanism when packets or
ACKs are lost.

\subsection{Stop and Wait}
The simplest configuration of the sender uses a window size of one, which
matches the stop and wait protocol (\textit{rdt 3.0}) in Kurose and Ross
\S3.4.1 \cite{b3}. The sender transmits a single packet, starts a timer, and
blocks for the matching ACK before sending the next packet. If the timer
expires first, the same packet is retransmitted. This case is the baseline for
all later experiments because every reliability mechanism (sequence numbers,
checksums, ACKs, and timeouts) is exercised without the additional bookkeeping
introduced by a larger window.

\subsection{Sliding Window}
With a window size greater than one, the sender can transmit multiple packets
before all acknowledgments arrive, which improves throughput on links with
non-trivial round-trip time. The implementation keeps track of unacknowledged
packets inside the active window and retransmits timed out packets
individually rather than rolling back the whole window. The receiver also
keeps a small buffer for valid packets that arrive inside its receive window,
so out of order packets can be acknowledged immediately and delivered upward
later when the missing earlier packets arrive.

\section{Implementation}

\subsection{Sender}
The sender handles the main reliability logic. It reads a message from the
command line or a file, splits the message into payload sized chunks, assigns
sequence numbers, computes checksums, and sends the packets over a UDP socket.
For each outstanding packet, it stores the packet data, the number of send
attempts, and the time the packet was last transmitted.

The sender accepts several parameters at runtime, including destination host and
port, timeout, payload size, window size, and maximum retries. This made it
possible to test both a stop and wait configuration and a larger sending window
without changing the code. The sender also records statistics such as original
packets, total packets sent, retransmissions, timeouts, acknowledgments
received, malformed acknowledgments, and checksum failures.

The sender parses each incoming ACK, checks whether it is malformed, and
validates its checksum. The receiver computes checksums for ACK packets, and the
sender also has an optional compatibility mode for older checksum zero ACKs.
The sender implementation covers sequence numbering, timeout recovery,
retransmission, and configurable window support.

Packet construction is done before any data is sent. Each chunk gets the next
sequence number, and the checksum is computed from the packet fields.

\begin{lstlisting}[language=Python]
packets = []
for offset, payload in enumerate(chunks):
    self._validate_payload(payload)
    seq_num = self.seq_num + offset
    checksum = calculate_checksum(seq_num, 0, payload)
    packets.append(Packet(seq_num, 0, payload, checksum))
return packets
\end{lstlisting}

The sender tracks unacknowledged packets by sequence number. On each loop, it
sends any unsent packets that still fit inside the current window.

\begin{lstlisting}[language=Python]
window_base = min(states)
window_limit = window_base + self.window_size

while next_to_send <= final_seq \
        and next_to_send < window_limit:
    if next_to_send in states:
        self._send_packet(states[next_to_send])
    next_to_send += 1

ack_num = self._receive_ack()
if ack_num in states:
    del states[ack_num]
\end{lstlisting}

ACKs are only accepted if they can be parsed and pass the sender's checksum
check.

\begin{lstlisting}[language=Python]
try:
    ack_packet = extract_from_bytes(ack_bytes)
except (UnicodeDecodeError, ValueError, IndexError):
    self.stats.malformed_acks += 1
    return None

if not self._ack_checksum_is_valid(ack_packet):
    self.stats.bad_ack_checksums += 1
    return None

return ack_packet.ack_num
\end{lstlisting}

Retransmission is also handled by sequence number. If a packet has not been
acknowledged before its timeout expires, the sender sends that packet again.

\begin{lstlisting}[language=Python]
if last_sent is not None and now - last_sent >= self.timeout:
    self.stats.timeouts += 1
    self._send_packet(states[seq_num])
\end{lstlisting}

\subsection{Receiver}
The receiver handles the validation and delivery side of the protocol. It binds
a UDP socket, waits for incoming packets, parses each packet, and checks the
packet checksum before accepting the data. The receiver keeps a receive window,
a buffer for valid packets, and the next expected sequence number. This allows
it to accept packets that arrive out of order while still delivering data upward
only in sequence order.

The receiver also builds ACK packets with the acknowledged sequence number. The
ACK checksum is computed over the ACK fields before the packet is sent back to
the sender.

\begin{lstlisting}[language=Python]
def send_ack(self, ack_num, addr):
    checksum = calculate_checksum(0, ack_num, '')
    p = Packet(0, ack_num, '', checksum)
    self.s.sendto(convert_to_bytes(p), addr)
\end{lstlisting}

When a packet arrives, the receiver recalculates the checksum from the received
sequence number, acknowledgment number, and payload. If the checksum does not
match, the packet is ignored and no payload is delivered.

\begin{lstlisting}[language=Python]
calculated_checksum = calculate_checksum(
    p.seq_num, p.ack_num, p.payload
)
if p.checksum != calculated_checksum:
    continue
\end{lstlisting}

After a packet passes the checksum check, the receiver compares the sequence
number to the current receive window. If the packet was already received
recently, the receiver sends the ACK again. This handles the symmetric loss
case: when the original ACK is dropped, the sender's timer fires and it
retransmits the same sequence number, so the receiver must re-acknowledge it
or the sender will keep retransmitting until \texttt{max\_retries} is reached
and the transfer aborts.

\begin{lstlisting}[language=Python]
if p.seq_num >= self.expected_seq_num - self.window_size \
        and p.seq_num < self.expected_seq_num:
    self.send_ack(p.seq_num, addr)
\end{lstlisting}

If the packet is inside the current receive window, the receiver acknowledges it
and stores it in the buffer. The receiver then slides forward while the next
expected sequence number is already present in the buffer.

\begin{lstlisting}[language=Python]
elif p.seq_num >= self.expected_seq_num \
        and p.seq_num < self.expected_seq_num + self.window_size:
    self.send_ack(p.seq_num, addr)

    if p.seq_num not in self.buffer:
        self.buffer[p.seq_num] = p

    while self.expected_seq_num in self.buffer:
        self.buffer.pop(self.expected_seq_num)
        self.expected_seq_num += 1
\end{lstlisting}

Packets outside the receive window are ignored. This keeps the receiver from
accepting packets that are too far ahead of the current delivery point or no
longer relevant to the active window.

\subsection{Testing Scripts}
Three Python scripts in \texttt{scripts/} drive the experiments, and a fourth
manual smoke test was used during development.

The smoke test runs the sender and receiver in two terminals without any
wrapper:

\begin{lstlisting}[language=bash]
python3 receiver.py
python3 sender.py --message "testing a larger sender window" \
    --payload-size 4 --window-size 10
\end{lstlisting}

This run confirmed that messages were split into the expected number of
packets, that ACKs were received in order, and that the receiver delivered
the data upward without any artificial impairment.

\texttt{run\_preset\_tests.py} starts a fresh receiver, runs the sender
against a named preset, and writes per-run sender and receiver logs. The
script also exposes an interactive \texttt{--menu} option that lists the
available presets later summarized in Table~\ref{tab:setup}, so that one
impairment at a time can be exercised during a demonstration.

\texttt{run\_proxy\_tests.py} inserts an in-process proxy between the sender
and the receiver. The proxy applies deterministic delay, jitter, loss,
corruption, duplication, and reordering according to a preset, so the
impairments are reproducible across runs without depending on
\texttt{tc netem} configuration on the host. The proxy reordering preset is
the test that most clearly distinguishes sender-side ordering from receiver
arrival order, because the sender log shows sequence numbers leaving in
order while the receiver log shows them arriving out of order before being
buffered and delivered in sequence.

\texttt{summarize\_results.py} parses the saved logs and reports the per-run
counters defined by \texttt{SenderStats}: original packets, packets sent,
retransmissions, timeouts, ACKs received, bad ACK checksums, malformed ACKs,
and unexpected ACKs.

\section{Experimental Setup}
Each preset fixes a network impairment profile and the sender configuration
used to run against it. Table~\ref{tab:setup} summarizes the presets used in
the local tests. Window size, timeout, payload size, and maximum retries are
the same sender parameters exposed on the command line, so the tests cover
both stop and wait (window size $1$) and a small sliding window (window size
$4$). The raw logs are grouped under
\texttt{linux\_test\_results/01\_two\_terminal\_manual\_demo/},
\texttt{linux\_test\_results/02\_tc\_netem\_preset\_tests/}, and
\texttt{linux\_test\_results/03\_proxy\_impairment\_tests/}. A final
intentional failure test is saved under
\texttt{linux\_test\_results/04\_limit\_failure\_test/}, and a high-loss
threshold test is saved under
\texttt{linux\_test\_results/05\_loss\_threshold\_test/}. A timeout tuning
test is saved under
\texttt{linux\_test\_results/06\_timeout\_threshold\_test/}. The proxy presets
(\texttt{all\_tests} and \texttt{rigorous\_test}) combine multiple impairments
at once and use a longer message so the sender has to split it across many
packets.

\begin{table}[t]
\caption{Test presets. ``--'' = default sender retry budget.}
\label{tab:setup}
\centering
\scriptsize
\setlength{\tabcolsep}{3pt}
\resizebox{\columnwidth}{!}{%
\begin{tabular}{|l|c|c|c|c|c|c|c|}
\hline
\textbf{Preset} & \textbf{Delay} & \textbf{Loss} & \textbf{Corrupt} & \textbf{Duplicate} & \textbf{Reorder} & \textbf{Window} & \textbf{Retries} \\
\hline
baseline & 0 & 0 & 0 & 0 & 0 & 1 & -- \\
delay\_100ms & 100\,ms & 0 & 0 & 0 & 0 & 1 & -- \\
loss\_5\_percent & 0 & 5\% & 0 & 0 & 0 & 1 & -- \\
corrupt\_2\_percent & 0 & 0 & 2\% & 0 & 0 & 1 & -- \\
corrupt\_10\_percent & 0 & 0 & 10\% & 0 & 0 & 1 & -- \\
window\_4\_baseline & 0 & 0 & 0 & 0 & 0 & 4 & -- \\
delay\_100ms\_loss\_5\_percent & 100\,ms & 5\% & 0 & 0 & 0 & 1 & -- \\
extreme\_test & 100$\pm$20\,ms & 10\% & 5\% & 2\% & 5\% & 4 & 50 \\
all\_tests (proxy) & 100$\pm$20\,ms & 5\% & 2\% & 1\% & 2\% & 4 & 40 \\
rigorous\_test (proxy) & 100$\pm$30\,ms & 10\% & 5\% & 2\% & 5\% & 1 & 80 \\
stress\_loss\_100 & 0 & 100\% & 0 & 0 & 0 & 1 & 3 \\
stress\_loss\_70 & 0 & 70\% & 0 & 0 & 0 & 1 & 3 \\
stress\_delay\_500\_timeout\_100ms & 500\,ms & 0 & 0 & 0 & 0 & 1 & 3 \\
\hline
\end{tabular}}
\end{table}

\section{Results}
Table~\ref{tab:results} reports the counters parsed from the sender logs in
\texttt{linux\_test\_results/} for each preset. Every non-stress preset listed
completed successfully, meaning the sender exited only after every packet was
acknowledged; the three stress entries failed by design.

\begin{table}[t]
\caption{Sender stats per preset (from \texttt{linux\_test\_results/}).}
\label{tab:results}
\centering
\scriptsize
\setlength{\tabcolsep}{3pt}
\resizebox{\columnwidth}{!}{%
\begin{tabular}{|l|c|c|c|c|c|c|}
\hline
\textbf{Preset} & \textbf{Result} & \textbf{Original} & \textbf{Sent} & \textbf{Retransmissions} & \textbf{Timeouts} & \textbf{ACKs} \\
\hline
baseline & Pass & 10 & 10 & 0 & 0 & 10 \\
delay\_100ms & Pass & 10 & 10 & 0 & 0 & 10 \\
loss\_5\_percent & Pass & 10 & 12 & 2 & 2 & 10 \\
corrupt\_2\_percent & Pass & 11 & 11 & 0 & 0 & 11 \\
corrupt\_10\_percent & Pass & 11 & 13 & 2 & 2 & 11 \\
window\_4\_baseline & Pass & 11 & 11 & 0 & 0 & 11 \\
delay\_100ms\_loss\_5\_percent & Pass & 8 & 9 & 1 & 1 & 8 \\
extreme\_test & Pass & 29 & 46 & 17 & 17 & 29 \\
all\_tests & Pass & 17 & 59 & 42 & 42 & 17 \\
rigorous\_test & Pass & 46 & 62 & 16 & 16 & 49 \\
stress\_loss\_100 & Fail & 8 & -- & -- & -- & 0 \\
stress\_loss\_70 & Fail & 9 & -- & -- & -- & 0 \\
stress\_delay\_500\_timeout\_100ms & Fail & 10 & -- & -- & -- & 0 \\
\hline
\end{tabular}}
\end{table}

In the clean cases (baseline, delay only, and the window-of-four baseline),
\texttt{packets\_sent} equals \texttt{original\_packets} and there are no
timeouts or retransmissions, which confirms the sender, receiver,
packetization, CRC, ACK, and delivery path on the happy path. The
\texttt{delay\_100ms} run shows that the default $1.0\,\text{s}$ timeout is
larger than the added round-trip delay, so the sender does not retransmit
unnecessarily.

The lossy and corrupt cases show the recovery path. With $5\%$ loss the
sender saw two timeouts and recovered with two retransmissions, ending with
$10$ ACKs for $10$ original packets. The $10\%$ corruption case is similar:
two corrupt packets are dropped by the receiver's CRC check, the sender's
timer fires, and two retransmissions complete the transfer. The $2\%$
corruption case finished with no retransmissions because no packet happened
to be corrupted in that particular run.

The combined \texttt{tc netem} presets show the same recovery behavior under
multiple impairments. \texttt{delay\_100ms\_loss\_5\_percent} needed one
retransmission to deliver eight original packets. \texttt{extreme\_test}
combined delay, jitter, loss, corruption, duplication, and reordering with a
window size of four; it still completed, but it required $46$ packets sent for
$29$ original packets.

The two proxy presets stress the protocol with several impairments at once.
\texttt{all\_tests} sent $59$ packets to deliver $17$ original packets, with
$42$ retransmissions and $42$ timeouts, all driven by the proxy's combined
loss, corruption, duplication, and reordering. \texttt{rigorous\_test} runs in
stop and wait over a longer message: $46$ original packets, $62$ packets sent,
$16$ retransmissions, and $49$ ACKs received. Three of those ACKs failed the
sender's CRC check (\texttt{bad\_ack\_checksums=3}) and were ignored, which is
the expected behavior when the proxy corrupts an ACK on the return path.

The intentional failure test used Linux \texttt{tc netem} with $100\%$ packet
loss, window size $1$, a small sub-second timeout, and
\texttt{max\_retries=3}.
The sender repeatedly retransmitted sequence number $0$ and then exited with
\texttt{TimeoutError: no valid ACK for seq=0 after 3 retransmission(s)}. The
receiver log contains only the startup message, which is consistent with the
network dropping every packet before it reached the receiver.

The threshold-style failure test used $70\%$ packet loss with the same retry
cap and a small sub-second timeout. In that run, the receiver did receive and ACK sequence number
$0$, but the sender still did not receive a valid ACK before the retry limit.
This shows that the failure point is not only complete loss; a high enough loss
rate can also exhaust the configured retry budget.

The more realistic timeout-threshold test used no packet loss, but added
$500\,\text{ms}$ of delay while the sender timeout was only
$100\,\text{ms}$. The receiver log shows that sequence number $0$ arrived and
was ACKed multiple times. However, the sender retransmitted and hit
\texttt{max\_retries=3} before a delayed ACK arrived back in time. This
demonstrates a practical limitation: even when packets are not lost, a timeout
that is too small for the path delay can make a reliable transfer fail.

\section{Discussion}
The results show that the implementation matches the textbook RDT model
inside its configured limits. The CRC catches corrupt packets in both
directions. Sequence numbers handled the duplicates the proxy injected,
and retransmission covered the packet losses. Out-of-order arrivals were
absorbed by the receiver's buffer so the upper layer still sees ordered
data. The window-of-four cases confirm that the sliding window logic does
not break correctness on the easy paths.

The protocol still has a clear configured limit even though the final
\texttt{extreme\_test} run completed successfully. The
\texttt{stress\_loss\_100}, \texttt{stress\_loss\_70}, and
\texttt{stress\_delay\_500\_timeout\_100ms} runs show that a single sequence
number that is repeatedly lost or repeatedly unacknowledged blocks progress.
Once its send attempt count exceeds
\texttt{max\_retries}, the sender raises \texttt{TimeoutError} and aborts
rather than hanging forever. This is the right failure mode for a fixed-retry
sender, but it also shows that under harsher conditions the protocol depends on
retry budget and timeout values being tuned to the network, not on the
mechanisms themselves.

A second observation comes from the \texttt{packets\_sent} versus
\texttt{original\_packets} ratio. In the clean cases the ratio is $1.0$; in
the proxy \texttt{all\_tests} run it is roughly $3.5$. That overhead is the
direct cost of reliability under combined loss, corruption, and reordering,
and it grows with the impairment rate as expected.

\section{Conclusion}
This project added reliability to a UDP-based transport. The sender uses
sequence numbers, CRC checksums, and timers with retransmission. The
receiver validates checksums, acknowledges valid packets, and buffers
in-window arrivals so it can deliver data in order. Both sides support a
configurable window so the protocol can run as stop and wait or with
multiple packets in flight. These are the same RDT mechanisms covered in
the textbook.

The same sender and receiver code passed clean baselines,
single-impairment tests, and combined proxy stress tests. When the network
got harsher than the sender was configured for, the protocol did not
deadlock; it hit the retry cap and exited with a clear error. The three
aggressive stress tests ($100\%$ loss, $70\%$ loss, and $500\,\text{ms}$
delay against a $100\,\text{ms}$ timeout) confirm this boundary. The
implementation is small enough to read end to end, which made it useful
for checking each RDT idea on its own.

\section{Future Work}
Several improvements would make the implementation closer to a production
transport. The text packet format should be replaced with a binary-safe
encoding so payloads can contain the \texttt{|} separator and arbitrary
file bytes. The receiver should be able to write delivered payloads to a
file rather than only printing them, and the protocol should add an
explicit end-of-message or connection teardown packet so the receiver
knows when a transfer is complete. An adaptive timeout based on measured
round-trip samples, similar to TCP's RTT estimation, would let the sender adjust
to path delay automatically rather than relying on a fixed value \cite{b3}. On the testing
side, automated unit tests and CI-style regression tests would help catch
protocol edge cases, and additional sweeps across window sizes and loss rates
would show how throughput trades off against reliability.

\section*{Acknowledgment}
The authors thank the CS 5470 instructor for course material on reliable
data transfer, and reference the textbook by Kurose and Ross for the RDT
progression and the CRC presentation that this project follows.

\begin{thebibliography}{00}
\bibitem{b1} W. Eddy, ``Transmission Control Protocol (TCP),'' RFC 9293,
Internet Engineering Task Force, Aug. 2022.

\bibitem{b2} J. Iyengar and M. Thomson, ``QUIC: A UDP-Based Multiplexed and
Secure Transport,'' RFC 9000, Internet Engineering Task Force, May 2021.

\bibitem{b3} J. F. Kurose and K. W. Ross, \textit{Computer Networking: A
Top-Down Approach}, 8th ed. Hoboken, NJ, USA: Pearson, 2021.

\end{thebibliography}

\end{document}
